\section{第五天}
\begin{enumerate}
  \item 核心要义之\bs{坚持依宪治国、依宪执政}。
        \begin{enumerate}
          \item 依宪治国、依执政的概念和价值。
                \begin{enumerate}
                  \item 一个关键句，两个展开逻辑。
                        \begin{enumerate}
                          \item 坚持依法治国\zx{首先要}坚持依宪治国，坚持依法执政\zx{首先要}坚持依宪执政。
                          \item 这是由宪法的\bs{根本大法地位和最高法律效力}所决定的。
                        \end{enumerate}
                  \item 根本大法地位 = 根本大法 + 地位。
                        \begin{enumerate}
                          \item 宪法是\zx{国家的根本大法}，是\zx{治国理政的总章程}，是\zx{党和人民意志的集中体现}。
                          \item 宪法\zd{规定了}国家的\zx{根本制度和根本任务}，\zd{保障了}\zx{公民的基本权利}，
                          \item 具有\bs{最高的法律权威和法律地位}。
                        \end{enumerate}
                  \item 最高法律效力。
                        \begin{enumerate}
                          \item 一切法律、行政法规、地方性法规\zd{均不得}同宪法相\bs{抵触}。
                          \item 一切国家机关、武装力量、社会团体和公民\zd{都要在}宪法和法律规定的\bs{范围内活动}。
                          \item 一切违反宪法和法律的行为都要\bs{予以追究}。
                        \end{enumerate}
                \end{enumerate}
          \item 如何（主要背第一句）。
                \begin{enumerate}
                  \item 全面贯彻\zx{实施宪法}。
                        \begin{enumerate}
                          \item 规定国家\zx{宪法日}，
                          \item 国家公职人员就职前\zx{宪法宣暂}。
                        \end{enumerate}
                  \item 推进\zx{合宪性审查}工作。
                        \begin{enumerate}
                          \item \zd{强化}全人大和全人常\zx{宪法监督}职责，
                          \item \zd{完善}全人常\zx{宪法解释}程序，
                          \item \zd{加强}\zx{备案审查制度和能力建设}。
                        \end{enumerate}
                  \item 深入开展\zx{宪法教育宣传}。
                        \begin{enumerate}
                          \item \zx{人民}群众\zx{认识宪法、认同宪法}，让宪法\zx{家喻户晓、深入人心}。
                          \item \zx{领导干部}加强\zx{宪法学习}，带头\zx{尊崇、学习、遵守和运用}宪法。
                        \end{enumerate}
                \end{enumerate}
        \end{enumerate}
  \item 核心要义之\bs{坚持统筹推进国内法治和涉外法治}。
        \begin{enumerate}
          \item 概念和价值。
                \begin{enumerate}
                  \item 统推内法外治是维护国家\bs{主权、安全、发展}利益的\zx{迫切需要}。\bs{内因}。
                  \item 当今世界正面临百年未有之\zx{大变局}，国际社会\zx{经济发展和地缘政治}安全发生深刻\zx{变化}。\bs{外因}。
                  \item 国家主权、安全、发展利益是国家的\bs{核心利益}，\bs{根因}。
                  \item \zx{切实维护}国家主权、安全、发展利益是涉外法治工作的\bs{首要任务}。
                  \item 当前，随着我国\zx{经济实力和综合国力}快速\zx{增长}，
                  \item \zx{对外开放}全方位\zx{深化}，\zx{“一带一路”建设}深入\zx{推进}，
                  \item 我国日益\zd{走近}世界\zx{舞台中央}，深度\zx{融入全球化}进程，\bs{宏观}。
                  \item \zd{维护}我国国家利益和\zx{公民、法人在境外}合法权益的任务\zx{日益繁重}。\bs{微观}。
                  \item 统推内法外治，协调推进\zx{国内治理和国际治理}，
                        \begin{enumerate}
                          \item 是\zx{全治国的必然要求}，
                          \item 是建立以\bs{国内大循环为主体、国内国际双循环相促进的新发展格局}的客观需要，
                          \item 是维护国家主权、安全、发展利益的\zx{迫切需要}。
                        \end{enumerate}
                \end{enumerate}
          \item 如何。
                \begin{enumerate}
                  \item 加快\zx{涉外法治工作布局}。\zd{完善制度}、加强\zd{理论研究}、加强\zd{涉外法治队伍}建设。
                  \item 加强\zx{国际法治合作}。积极参与国际\zx{新秩序构建}、发展对外\zx{友好关系}、加强\zx{反腐}国际合作和国际\zx{司法协助}。
                  \item 为构建\zx{人类命运共同体}提供法治保障。
                        \begin{enumerate}
                          \item 维护\zx{以联合国为核心}的国际体系。
                          \item 为\zx{人类进步事业}贡献\zx{中国方案、中国智慧}。
                        \end{enumerate}
                \end{enumerate}
        \end{enumerate}
\end{enumerate}